\clearpage
\fancyhead[L]{\small INSTRUMENTS.}
\hypertarget{native-429-part-lower}{}
\begin{center}{\Large INSTRUMENTS.}\end{center}
\begin{center}
Astronomisches Jahrbuch für 1822. pp. 235..241. Berlin 1819.
\end{center}

\par\indent I first observed the comet on 2 July in Lauenburg, on the evening when Lieutenant \textsc{Zartmann} found it there. I did not, however, reduce these observations precisely (they were made by measuring distances from stars with reflecting sextants), since I soon learned that around the same time the comet had already been discovered elsewhere and observed on the meridian. After my return, I observed it here on the meridian at its lower culmination, initially through 26 July with \textsc{Repsold}’s circle, where, however, the crossings of the fine threads began to be difficult to observe, while the zenith distances still to allow greater precision. 
In the meantime I had made \textsc{Reichenbach}’s transit telescope ready for observations. Along with its considerably greater optical power, this telescope has somewhat thicker threads. So from 27 July onward I observed the right ascensions with this instrument, while Professor \textsc{Harding} continued to set the zenith distance on \textsc{Repsold}’s circle. Since 4 August I have seen the comet through the transit telescope on several more evenings (for the last time on 21 August). But since the \SourceLeaf{430}thread crossings were guessed at more than actually observed, I prefer to omit them as unworthy of confidence. They do, however, confirm the continued good agreement of \textsc{Dircksen}’s elements.

\Needspace{18\baselineskip}
\begin{center}
\emph{Observations of the comet.}
\end{center}

\begin{center}
\begin{tabular}{l|c|c|c}
1819 & Mean time & & \\
\hline
July 19 & $11^{\mathrm u}\,56^{\mathrm m}\,53^{\mathrm s}$ & $116^\circ\,7'\,26''$ & $+51^\circ\,52'\,12''$ \\
21 & $11\,53\,33$ & $117\,15\,34$ & $51\,54\,58$ \\
22 & $11\,51\,53$ & $117\,49\,34$ & $51\,54\,56$ \\
26 & $11\,44\,14$ & $119\,51\,10$ & $51\,48\,49$ \\
27 & $11\,42\,8$ & $120\,18\,43$ & --- --- --- \\
28 & $11\,40\,2$ & --- --- --- & $51\,43\,46$ \\
30 & $11\,35\,40$ & $121\,38\,52$ & $51\,37\,23$ \\
Aug. 3 & $11\,26\,24$ & $123\,16\,9$ & $51\,23\,25$ \\
4 & $11\,23\,59$ & $123\,38\,55$ & $51\,19\,46$ \\
\end{tabular}
\end{center}

\Needspace{13\baselineskip}
\par\indent Mr. \textsc{Dircksen} has calculated the following parabolic elements from these observations and the first observations made by Professor \textsc{Struve} in Dorpat.

\begin{center}
\begin{tabular}{ll}
Time of perihelion & 1819 June 28.\ $17^{\mathrm n}\,49^{\mathrm m}\,9^{\mathrm s}$ Göttingen time \\
Longitude of perihelion & . . . . . . . . . . .\ $287^\circ\,6'\,24''8$ \\
Logarithm of the distance & . . . . . . . . . .\ $9.5330800$ \\
Longitude of the node & . . . . . . . . . . . .\ $273^\circ\,42'\,8''8$ \\
Inclination of the orbit & . . . . . . . . . . . .\ $80\,45\,12.0$ \\
\multicolumn{2}{l}{Motion prograde.}
\end{tabular}
\end{center}

\Needspace{17\baselineskip}
\par\indent Comparing the observations above with these elements gives the following differences:

\begin{center}
\begin{tabular}{l|c|c}
 &  Right asc.  &  Decl.  \\
\hline
July 19  &  $-15''5$  &  $-3''5$  \\
21  &  $+39.7$  &  $+5.2$  \\
22  &  $-30.0$  &  $+9.5$  \\
26  &  $-37.9$  &  $+26.5$  \\
27  &  $-9.8$  &  --- ---  \\
\hline
  &  Right asc.  &  Decl.  \\
\hline
 July 28  &  --- ---  &  $+4''0$  \\
 30  &  $-0''2$  &  $-19.2$  \\
 Aug. 3  &  $-9.1$  &  $-0.3$  \\
 4  &  $+3.8$  &  $-2.9$  \\
\end{tabular}
\end{center}

\Needspace{18\baselineskip}
\par\indent Comparing some earlier observations made elsewhere gives the following differences:

\begin{center}
\begin{tabular}{l|c|c|l}
July 2  &  $-25''1$  &  $-16''6$  &  Bode  \\
3  &  $-1.0$  &  $-1.6$  &  Struve  \\
3  &  $-1.8$  &  $+44.4$  &  Pond  \\
4  &  $-17.3$  &  $+19.6$  &  Struve  \\
4  &  $-48.3$  &  $-65.8$  &  Bode  \\
5  &  $+12.1$  &  $+18.8$  &  Struve  \\
\hline
 July 5  &  $-19''1$  &  $-86''7$  &  Pond  \\
 6  &  $-16.6$  &  $+18.0$  &  Struve  \\
 6  &  $-20.8$  &  $-19.9$  &  Bode  \\
 7  &  $-18.1$  &  $+17.3$  &  Struve  \\
 11  &  $+41.1$  &  $-24.4$  &  Bode  \\
 12  &  $-96.3$  &  $-6.3$  &  Bode  \\
\end{tabular}
\end{center}

\par\indent The six-foot transit telescope made by \textsc{Reichenbach} has now been in daily use since the end of July. I think you will welcome a brief account of this instrument. The objective by \textsc{Frauenhofer} has an aperture of 52 Paris lines. Its 4 eyepieces magnify (according to my still preliminary estimate) 80, 100, 140, 200 times. I usually work at the highest magnification, though careful setting is needed to avoid missing the stars, since the radius of the field of view is only $4\tfrac{1}{2}$ minute. There are 7 vertical threads, with intervals that stars on the equator traverse within $10^{\mathrm s}$. With a little practice one readily gets used to these short intervals. In the time between crossings it is quite possible to write down the observations, slide the eyepiece, \SourceLeaf{431}and finish early enough to wait with complete composure for the two beats of the clock between which the next crossing falls. One only has to be accustomed to keeping an accurate count of the beats throughout these operations, if one cannot have an assistant count them. Once this skill has been acquired, one finds it extremely agreeable to obtain a greater number of determinations in a short time. My own experience already shows that under that condition the greater number and closer spacing of the threads do not detract from the accuracy of the individual observations. From nearly 300 observed crossings of $\gamma$, $\alpha$ and $\beta$ in Aquila I find the probable error of \emph{one} crossing to be $= 0^{\mathrm{s}}090$, a little smaller still than in my observations with \textsc{Repsold}’s circle, where the 5 threads are over $16^{\mathrm{s}}$ apart in time. — The telescope is set by means of a small circle, fixed to the side of the tube near the eyepiece and divided into quarter degrees, on which an alidade with a level moves. The vernier on this alidade indicates single minutes. This arrangement is similar to the one on the new Greenwich transit telescope. Its great advantage is that one can set the instrument just as easily when it is reversed. Once the alidade has been set, one can also return the tube to the correct position more quickly if it has accidentally shifted a little. The setting, however, seems to take a little more time than with the usual arrangement, since it requires \emph{two} operations. At night one also has difficulty managing without an assistant. Mr. \textsc{von Reichenbach} will supply me with a second circle of a somewhat different design, which will remove this last inconvenience and have the further advantage that one can set the two circles in advance for two stars and thus readily observe them even if they follow very quickly one after the other. I also give you my observations of the last opposition of Jupiter, for which I observed the right ascensions with the transit telescope while Professor \textsc{Harding} set the zenith distances on \textsc{Repsold}’s circle. These observations cannot yet serve as a measure of the accuracy attainable with the transit telescope, because the clock used at the time showed some irregularities in its rate.

\begin{center}
\begin{tabular}{r|ccc|cc|cc}
1819 & \multicolumn{3}{c|}{Mean time} & \multicolumn{2}{c|}{} & \multicolumn{2}{c}{} \\
\hline
July 30 & $12^{\mathrm{u}}$ & $31^{\mathrm{m}}$ & $46^{\mathrm{s}}4$ & $315^{\circ}\,42'$ & $47''3$ & $-17^{\circ}\,47'$ & $50''8$ \\
Aug.\ 1 & 12 & 22 & 52.5 & $27$ & 13.2 & $52$ & 34.8 \\
3 & 12 & 13 & 58.7 & $11$ & 41.1 & $57$ & 12.1 \\
4 & 12 & 9 & 31.7 & $3$ & 53.5 & $59$ & 35.9 \\
5 & 12 & 5 & 4.5 & $314\,56$ & 3.1 & $18\,1$ & 49.3 \\
13 & 11 & 29 & 29.1 & $313\,53$ & 50.7 & $19$ & 56.1 \\
\end{tabular}
\end{center}

\par\indent Last autumn I received from Mr. \textsc{Repsold} an excellent new level which is placed directly on the telescope. Comparing observations made with the instrument in opposite positions gives the angle between the leveled line and the line of sight. Since then I have made the observation of about 50 circumpolar stars one of my main occupations and have obtained about 1200 observations. I will publish the results for the altitude of the pole and the declinations of these stars on another occasion, but several other tests must come first. For one thing I will first examine the graduations as carefully as possible with the aid of a special apparatus now fitted to the piers. For another I will make a further series of observations which I consider indispensable if the \emph{absolute} zenith distances are to be reliable. Much recent experience with stellar declinations and the obliquity of the ecliptic has shown that we are \SourceLeaf{432}still quite uncertain about the true absolute zenith distances. One instrument gives one result and another gives a different one. Since we cannot eliminate the effect of \emph{gravity} on the parts of the instrument, flexure and the like give rise to \emph{errors}, depending on the construction of the instrument. They vary with zenith distance and may be assumed to have the form $a\sin z+b\cos z$. In instruments that can be reversed and with which every star is observed in both positions, the second term $b\cos z$ is harmless. It merely produces a collimation error that varies with the zenith distance, though even in a well-made instrument such a term will still be noticeable. The first term $a\sin z$ is a different matter. It distorts all zenith distances, the small ones a little and the larger ones more. In this respect each instrument is an individual. Although the instrument maker can greatly reduce this error, I consider it impossible to remove it completely with certainty, and the only course is to determine it by observation. Comparing the declinations of southern stars that Mr. \textsc{Littrow} derived from \textsc{Bessel}’s observations with those of \textsc{Pond} shows me that the difference between the values of $a$ for \textsc{Pond}’s and \textsc{Bessel}’s circles must be set at about $=3''$ to bring the results into the best agreement, but the individual values remain undetermined. If one knew \emph{for certain} that the center of the figure of the Sun describes a great circle and that refraction is exactly the same for sunlight and starlight, one would be entitled to attribute that difference wholly or for the most part to \textsc{Pond}’s circle. Since those assumptions are unproved, however, that conclusion would be a petitio principii. I will add only that comparing my observations with \textsc{Pond}’s shows that the value of $a$ for \textsc{Repsold}’s circle differs from that for \textsc{Pond}’s by $1''5$, on the opposite side from \textsc{Bessel}’s circle, so that the latter differs from \textsc{Repsold}’s by $4''5$. This should cause no surprise. On several occasions a weight of 1 Loth placed near the objective of \textsc{Repsold}’s circle produced flexure that reduced the zenith distances at the horizon by $1''$. The counterbalances fitted to the telescope by \textsc{Repsold} and \textsc{Reichenbach} do make this flexure very small, but cannot guarantee that it has been eliminated altogether. The only way to find the value of $a$, it seems to me, is to make a sufficiently large number of observations using a mercury horizon. So far I have made observations of this kind only a few times, and on the Pole Star. They give a mean of $a=0$ almost to a hair’s breadth, but the observations would need to be much more numerous before I could be satisfied with that. Even so, I can already believe that the value of $a$ for \textsc{Repsold}’s circle could be nearly $5''$. Observations of this kind present several practical difficulties, especially with a large instrument.

\par\indent From my observations with \textsc{Repsold}’s circle I also give here the right ascensions of the Pole Star that I determined last May. The decimal fraction of the day is reckoned from the lower culmination on the date. The comparison is made with \textsc{Bessel}’s \emph{new} tables\footnote{Omitting the two small equations, however.} (the differences in the observations in last year’s yearbook were referred to \textsc{Bessel}’s older tables). For the column headed Weight, I refer you to last year’s yearbook (p. 214 [p. 415 of this volume]).

\Needspace{22\baselineskip}
\SourceLeaf{433}
\begin{center}
\begin{tabular}{r|l|l|l}
1819  &  \shortstack{Apparent\\R.A.}  &  \shortstack{Corrected diff. from\\Bessel’s tables}  &  Weight  \\
\hline
 &  $0^{\mathrm u}\,55^{\mathrm m}$  &   &   \\
May 1.\,0  &  $51^{\mathrm s}03$  &  $+0^{\mathrm s}07$  &  5.52  \\
2.75  &  51.27  &  $-$0.40  &  3.43  \\
3.25  &  52.39  &  $+$0.51  &  4.44  \\
6.00  &  52.97  &  $-$0.09  &  3.56  \\
7.25  &  54.33  &  $+$0.72  &  3.75  \\
\hline
 1819  &  \shortstack{Apparent\\R.A.}  &  \shortstack{Corrected diff. from\\Bessel’s tables}  &  Weight  \\
\hline
  &  $0^{\mathrm u}\,55^{\mathrm m}$  &   &   \\
 May 8.25  &  $54^{\mathrm s}83$  &  $+0^{\mathrm s}76$  &  3.43  \\
 9.75  &  54.95  &  $+$0.17  &  3.00  \\
 10.75  &  53.29  &  $-$1.98  &  3.40  \\
 16.50  &  57.21  &  $-$1.06  &  3.00  \\
 21.25  &  60.96  &  $+$0.03  &  2.67  \\
\end{tabular}
\end{center}

The mean correction on this basis would be only $= -0^{\mathrm s}03$.

\par\indent The meridian circle made by \textsc{Reichenbach} has also arrived recently and can soon be installed. As for the transit telescope, I would add that its optical power is extraordinarily great. I see the little star $5'$ from the pole at twilight, when the threads still need no illumination. The instrument is reversed in the air by means of a device beautifully made by our local inspector, Mr. \textsc{Rumpf}. A rope runs over a pulley above, across to a second pulley at the side, and then down to the shaft of a winch with a ratchet, etc., around which it is wound. With it one raises the instrument only about 1 inch. It is then turned around while freely suspended and lowered back into place in the opposite position. After use the pulley is slid sideways simply by pulling a cord, so that it does not obstruct observations at the zenith. The whole operation takes about 5--6 minutes, so that if I reverse the instrument during the culmination of the Pole Star I miss only one of the 7 threads. Once I have received the second setting circle, it will be easy to carry out the operation so quickly that no thread is missed at all. With \textsc{Repsold}’s circle, reversal always takes over half an hour. With \textsc{Reichenbach}’s circle it will be done in a different way, by means of a special machine, and the instrument maker believes that 8 minutes will suffice each time.
